% Generated draft: requires mathematical and source review.
\documentclass[11pt]{article}
\usepackage[margin=25mm]{geometry}
\usepackage{fontspec}
\setmainfont{Latin Modern Roman}
\usepackage{amsmath,amssymb,mathtools,unicode-math}
\setmathfont{Latin Modern Math}
\usepackage{enumitem,needspace,xurl,hyperref}
\setlength{\parindent}{0pt}
\setlength{\parskip}{5pt}
\setlength{\emergencystretch}{3em}
\newcommand{\N}{\mathbb{N}}
\newcommand{\Z}{\mathbb{Z}}
\newcommand{\Q}{\mathbb{Q}}
\newcommand{\R}{\mathbb{R}}
\newcommand{\C}{\mathbb{C}}
\newcommand{\F}{\mathbb{F}}
\newcommand{\A}{\mathbb{A}}
\newcommand{\PP}{\mathbb{P}}
\newcommand{\E}{\mathbb{E}}
\newcommand{\eps}{\varepsilon}
\newcommand{\dd}{\,\mathrm{d}}
\begin{document}
\section*{UnsolvedMath ID: 800602}
\textbf{Source problem code:} catalogue entry
\subsection{Definitions and mathematical statement}
The supplied record does not specify the ambient class of spaces, admissible covers, or the precise splitting convention. The following is a conditional reconstruction using standard cover terminology; these conventions require comparison with the cited source.

Let \(\mathcal D\) denote the intended class of topological spaces. For each \(X\in\mathcal D\), let \(\mathcal C(X)\) be the intended collection of open covers of \(X\) whose members are proper subsets of \(X\). Whether these covers must be countable is unresolved in the supplied excerpt. Define
\[
\begin{aligned}
\Lambda(X)&=\{\mathcal U\in\mathcal C(X):\text{each }x\in X\text{ belongs to infinitely many }U\in\mathcal U\},\\
\Omega(X)&=\{\mathcal U\in\mathcal C(X):\text{each finite }F\subseteq X\text{ is contained in some }U\in\mathcal U\},\\
\Gamma(X)&=\{\mathcal U\in\mathcal C(X):\mathcal U\text{ is infinite and each }x\in X\text{ belongs to all but finitely many }U\in\mathcal U\}.
\end{aligned}
\]
A cover \(\mathcal U\in\Lambda(X)\) belongs to \(\mathrm T(X)\) if, for every \(x,y\in X\), at least one of the sets
\[
\{U\in\mathcal U:x\in U,\ y\notin U\},\qquad
\{U\in\mathcal U:y\in U,\ x\notin U\}
\]
is finite.

For this reconstruction, \(X\models\mathsf{Split}(\mathcal A,\mathcal B)\) means
\[
\forall\mathcal U\in\mathcal A(X)\ \exists\mathcal V_0,\mathcal V_1\subseteq\mathcal U:
\quad \mathcal V_0\cap\mathcal V_1=\varnothing,
\qquad \mathcal V_0,\mathcal V_1\in\mathcal B(X).
\]
Here the two subcovers need not exhaust \(\mathcal U\). Equivalently, \(\mathcal U\) can be partitioned into two families each containing a member of \(\mathcal B(X)\). The excerpt does not establish whether this convention, or the potentially stronger requirement that both parts themselves belong to \(\mathcal B(X)\), is intended.

With these qualifications, Problem 30 asks whether
\[
\forall X\in\mathcal D,\qquad
X\models\mathsf{Split}(\Omega,\mathrm T)
\ \Longrightarrow\
X\models\mathsf{Split}(\mathrm T,\Gamma).
\tag{1}
\]
If (1) is false, it asks whether nevertheless
\[
\forall X\in\mathcal D,\qquad
X\models\mathsf{Split}(\Omega,\mathrm T)
\ \Longrightarrow\
X\models\mathsf{Split}(\Lambda,\Lambda).
\tag{2}
\]
The intended choices of \(\mathcal D\), \(\mathcal C(X)\), and splitting convention must be recovered before this becomes a fully specified formulation of the source problem.
\subsection{Short English statement}
Does the property of splitting every \(\omega\)-cover into two disjoint \(\tau\)-subcovers imply the corresponding splitting property from \(\tau\)-covers to \(\gamma\)-covers? If that implication fails, does it still imply that every large cover has two disjoint large subcovers? This reading is conditional on unresolved ambient-space and cover conventions.
\subsection{Sources}
\begin{itemize}
\item Local UnsolvedMath registry, permanent ID \(800602\), exact source problem code \texttt{catalogue entry}, supplied complete source record. The registry attributes the problem to Tsaban, gives proposed year 2006, and labels it open; this status is supplied metadata, not an independently verified mathematical assessment.
\item Tsaban, \emph{Selection Principles and Special Sets of Reals: Open Problems} (2006), Problem 30. Supplied source URL: \url{https://arxiv.org/abs/math/0609601}. The record reports access on 2026-10-16 and extraction by \texttt{arxiv\_latex}; these are supplied provenance claims, not independent source verification.
\item The supplied attribution/context field contains \texttt{[split]}. No bibliographic expansion of this reference label is supplied, so its intended reference and any additional attribution remain unresolved.
\end{itemize}
\end{document}
